\subsection{Entrada, elección de constantes y salida por bandas}\label{meso:constantes}
\begin{lema}[Entrada regular en $M$]\label{meso:entrada}
Sean dadas las tolerancias de entrada de los observables de la Sección~\ref{meso:regularidad} y
$0<\alpha<1/4$. Supongamos que $\delta$ es tan pequeño que los errores deterministas de entrada, es
decir, las diferencias $O(\delta+1/n)$ entre los objetivos del tramo macroscópico en $M$ y los de esa
sección, caben con desigualdad estricta en esas tolerancias, uniformemente en las líneas distinguidas,
y que vale la segunda estimación del Lema~\ref{lema:presupuesto-meso} con una tasa $\lambda$. Si el
radio $\varrho_{\rm mac}$ y el paso $a_{\rm mac}$ de la Proposición~\ref{macro:parada} son
suficientemente pequeños y $n$ es suficientemente grande, en todo tablero de tamaño $M$ cuya
discrepancia $\Delta$ es menor que $\varrho_{\rm mac}$ los observables están dentro de sus
tolerancias de entrada, simultáneamente para todas las líneas distinguidas disponibles. Además, la
probabilidad de llegar a $M$ sin salida previa y con $D_{\max}(M)\ge\alpha M$ cumple
$\nu(G_M^-,D_{\max}(M)\ge\alpha M)\le C_\lambda e^{-\lambda M}$.
\end{lema}
\begin{proof}
Pongamos $x=M/n$. Uniformemente en $f$, se tiene $U_f=u_f+O(\delta+1/n)$ y $u_f\ge\kappa_\nu$, de
modo que $W_0U_DU_A=\nu_1(1+O(\delta+1/n))$. Para una fila o columna $\ell$ que sobrevive, sea
$F_\ell(M)$ la suma de los pesos $W_0$ de sus casillas disponibles. El objetivo macroscópico sobre
$\ell$ tiene densidad $xU_DU_A$ en la coordenada transversal. Al dividir por $x$ e integrar contra
$W_0$ se obtiene $1+O(\delta+1/n)$, pues la marginal de $\nu_1$ vale uno. En particular,
\[
 \left|\frac{F_\ell(M)}M-1\right|
 \le \frac{\|W_0\|_\infty}{x}\Delta_F+2\omega_F(a_{\rm mac})+C(\delta+1/n).
\]
Aquí $\omega_F(a_{\rm mac})$ acota la oscilación de $W_0$ en cada subfibra. Sustituir el peso por
un valor elegido en cada intervalo, tanto en la suma como en la integral objetivo, cuesta a lo sumo
el doble de esta oscilación en unidades residuales, pues cada masa total en unidades macroscópicas
es a lo sumo $x$.

Para un rectángulo de tipo $S\subseteq T$, sea
$w_S=\prod_{e\in S}W_{0,e}\prod_{e\in T\setminus S}\nu_{1,e}$, sea $Z_S(M)$ su observable inyectivo
y sea $I_*=\int_{[0,1]^4}\prod_{e\in T}\nu_{1,e}$ el objetivo
común. La densidad objetivo macroscópica para este tipo es $x^4\prod_{e\in S}U_{D(e)}U_{A(e)}$.
Tras dividir por $x^4$ y multiplicar por $w_S$, su integral es $I_*+O(\delta+1/n)$. Denotemos por
$\omega_S(a_{\rm mac})$ la oscilación máxima de $w_S$ en una caja de la partición. Entonces
\begin{align*}
 |Z_S(M)-I_*|
 &\le \frac{\|w_S\|_\infty}{x^4}\Delta_R+2\omega_S(a_{\rm mac})+C(\delta+1/n)\\
 &\quad +\frac{C}{M}+\|w_S\|_\infty\left|\frac{M^4}{(M)_2^2}-1\right| .
\end{align*}
En efecto, tanto la masa empírica ordinaria como la masa objetivo en unidades de $n^4$ suman a lo
sumo $x^4$ cada una. El término $C/M$ reincorpora las cuádruplas inyectivas con etiquetas diagonales
coincidentes. Cada coincidencia fija un índice una vez dados los otros tres, por lo que a lo sumo
$O(M^3)$ asignaciones de este tipo usan índices supervivientes. El último término cambia la
normalización de $M^4$ a $(M)_2^2$ y también es $O(1/M)$.

Elegimos la partición macroscópica de modo que refine la mesoscópica. En cada intervalo mesoscópico,
el error en la proporción residual de filas o de columnas es entonces a lo sumo $x^{-1}\Delta_0$.
Por la regularidad de Lipschitz, todas las oscilaciones anteriores son $O(a_{\rm mac})$,
uniformemente en las líneas distinguidas y en los dieciséis tipos. Una vez que $\delta$ deja un
margen estricto en las tolerancias de entrada, $x\ge\delta$ acota ambos factores de amplificación.
Podemos elegir $\varrho_{\rm mac}$, después un refinamiento de paso suficientemente fino como en la
Proposición~\ref{macro:parada} y finalmente $n$ suficientemente grande para los términos de redondeo
y de discretización. Con estas elecciones, las tres cotas quedan dentro de las tolerancias de entrada
siempre que $\Delta<\varrho_{\rm mac}$.

La cota para $D_{\max}(M)$ es la segunda estimación del Lema~\ref{lema:presupuesto-meso} en $M$.
Es una condición sobre todos los cruces vacíos y por ello se estima desde el tablero inicial, también
para los no disponibles.
\end{proof}
Decimos que una historia tiene \emph{entrada regular} en $M$ si no tiene salida hasta el tablero
de tamaño $M$, incluido éste. Un fallo de entrada en $M$, sea de un observable o de la geometría, es
una salida
macroscópica en el tablero de tamaño $M$ (Definición~\ref{def:regular}). Ese tablero dispone de las
cotas exteriores macroscópicas por la herencia desde su último padre regular. En particular el fallo
geométrico en $M$ tiene probabilidad a lo sumo $C_\lambda e^{-\lambda M}$, exponencial en $n$ para el
$\delta$ ya fijado.

Fijamos las constantes en el siguiente orden.
\begin{enumerate}
\item Partimos del perfil y de dos clases de la forma de~\cite[Sección~8.1]{pred}. La \emph{clase
interior} es la de las constantes $C_{\rm in},\eta_{\rm in},\gamma_{\rm in}$ de la
Proposición~\ref{meso:reconstruccion}, con
masa diagonal a lo sumo $\varepsilon$. La \emph{clase exterior}, que la contiene, tiene cotas de
difusión, vacancia y brecha con constantes $C_{\rm out},\eta_{\rm out},\gamma_{\rm out}$
estrictamente más débiles y masa diagonal a lo sumo $\varepsilon_{\rm out}=2\varepsilon$. El valor de
$\varepsilon$ se elige en el punto~4. Con las
cotas de la clase exterior, las constantes de las prescripciones, de los incrementos y de la herencia
por hijo se fijan una sola vez.
\item Elegimos $a_{\rm mes},\varrho_{\rm mes}$ y las tolerancias de las fibras, de los intervalos y
de cada uno de los dieciséis términos para que valgan el reequilibrado y la reconstrucción
espectral. Reservamos, por ejemplo, un cuarto de cada tolerancia individual para la entrada: son las
\emph{tolerancias de entrada}.
\item Fijamos $A\ge4C_{\rm geo}$. Con la familia de observables fijada, escogemos $\chi,b_1>0$ para
que $C_H\chi+b_1$ quepa estrictamente en otro cuarto de cada tolerancia, absorbiendo la constante
del presupuesto. La primera estimación del Lema~\ref{lema:presupuesto-meso} fija ahora $c_\chi>0$ y
las estimaciones de martingalas de la Proposición~\ref{prop:concentracion-meso} fijan tasas de
concentración positivas. Elegimos $c_{\rm pat}>0$ menor que todas estas tasas. Los parámetros de que
dependen ya se han fijado, independientemente de $\varepsilon$, de $\delta$ y del suelo.
\item Tomamos ahora $\varepsilon>0$ suficientemente pequeña para la reconstrucción, con
$\varepsilon_{\rm out}=2\varepsilon$ no mayor que la cota
$\varepsilon_0(C_{\rm out},\eta_{\rm out},\gamma_{\rm out})$ del Teorema~8.1 de~\cite{pred}, y para
que se cumpla $2\beta(\varepsilon_{\rm out})<c_{\rm pat}/4$, donde ${\beta(v)=2v/(1-v)}$. Fijamos
$\alpha=\varepsilon/(4\|W_0\|_\infty)<1/4$ y elegimos $\delta>0$ de modo que los errores
deterministas de entrada y el término $C_HA\delta$ quepan con desigualdad estricta en sus
tolerancias, la restricción $A(\delta+1/n)+\chi+1/K<1/4$ se cumpla para $n,K$ grandes y la
\emph{segunda} estimación del Lema~\ref{lema:presupuesto-meso} para el suceso geométrico
$D_{\max}(k)\ge\alpha k$ valga con algún $\lambda>c_{\rm pat}$.
\item Fijamos finalmente $K_0\ge k_0(C_{\rm out},\eta_{\rm out},\gamma_{\rm out})$, con el $k_0$ del
mismo teorema, y suficientemente grande; según el Lema~\ref{meso:entrada}, el radio
$\varrho_{\rm mac}$ y el paso $a_{\rm mac}$ del tramo macroscópico; y después $n$ suficientemente
grande para este $\delta$. Se absorben así los errores de redondeo, los términos $O(1/K_0)$ y los
restos macroscópicos. Las condiciones que definen la clase interior, incluida la masa diagonal
$\le\varepsilon$, son independientes del suelo $K$.
\end{enumerate}

\begin{proposicion}[Salida por bandas]\label{prop:salida-banda-meso}
Existen $C<\infty$ y $c>2\beta(\varepsilon_{\rm out})$ tales que para cada $K_0\le K\le M$ y $j\ge0$
con $2^jK\le M$,
\[
 \nu\big(\text{entrada regular en $M$ y primera salida en un tamaño }
              k\in[2^jK,2^{j+1}K)\big)
 \le C e^{-c\,2^jK}.
\]
\end{proposicion}
\begin{proof}
Para cualquier suelo $K'$ entre $K_0$ y $M$, aplicamos la Proposición~\ref{prop:concentracion-meso} y
la segunda estimación del Lema~\ref{lema:presupuesto-meso} con ese suelo. En una historia de entrada
regular, salvo sus excepciones, cada observable al llegar a cualquier tamaño $k\ge K'$ hasta la
primera salida queda dentro de su tolerancia de reconstrucción, y $D_{\max}(k)<\alpha k$. Si se
produjera la primera salida en este rango, su propio tablero de llegada cumpliría esas condiciones y
la Proposición~\ref{meso:reconstruccion} lo situaría en la clase interior. Por tanto la salida exige
una de las excepciones, de probabilidad $\le Ce^{-cK'}$ con $c=c_{\rm pat}/2$. Se toma $K'=2^jK$,
con la misma clase y las mismas constantes, para obtener el enunciado. El factor de coste de la
Sección~\ref{balance} es a lo sumo $\exp(C_0+2\beta(\varepsilon_{\rm out})K')$ en esta banda, por
lo que $c>2\beta(\varepsilon_{\rm out})$ deja una exponencial sumable en $K'=2^jK$. Tanto la
probabilidad como el presupuesto diagonal se toman desde el tablero inicial bajo $\nu$. Un paso que
llegue a $K$ y viole una condición se contabiliza como salida en $K$.
\end{proof}
